\chapter{Memoria descriptiva}
\thispagestyle{memoria}\pagestyle{memoria}

\section{Objeto}
El presente documento desarrolla el anteproyecto academico de instalaciones
electricas en baja tension para una estacion de servicio (grifo) de
combustibles liquidos. El trabajo es nuevo y pertenece a Renzo Gabriel Mamani
Galindo. Los planos DWG de ubicacion y distribucion recibidos se usan
exclusivamente para rescatar la implantacion, los ambientes, las islas, los
surtidores, la zona de tanques y los equipos electricos observados; no se
presentan como un proyecto electrico previo ni se modifica su fuente original.

El propietario del predio es Miguel Mamani Chuquicallata. La direccion del
predio es Rustico Reumita, Parcela B-8 y B-9 Lado Este, en la comunidad
campesina San Francisco de Buenavista, sobre la carretera Juliaca--Puno,
distrito y provincia de San Roman, departamento de Puno.

\section{Ubicacion y arquitectura}
El establecimiento ocupa un lote de aproximadamente 41~m~\(\times\)~37~m (lote
total) con un area a ejecutar de aproximadamente 28~m~\(\times\)~16~m. La
geometria se extrajo del plano de distribucion en coordenadas UTM WGS 84 zona
19S y se trabajo en coordenadas locales. La altitud del sitio es
aproximadamente 3830~m s.n.m.

La distribucion incluye un area administrativa (administracion y facturacion,
oficina 01), sala de maquinas, dos servicios higienicos, vereda y tuberias de
venteo. Se observan tres tanques enterrados (TK-1 Gasohol Premium, TK-2
Gasohol Regular, TK-3 Diesel B5-S50), dos islas de despacho, circulacion de
vehiculos mayores y menores, fosa de agua, cilindros de arena y trapo
empapado, y totem. GLP y GNV quedan expresamente fuera del alcance.

\section{Alcance electrico}
El anteproyecto incluye suministro de diseno 380/220~V, tablero general,
subalimentadores TDE (emergencia), TDF (fuerza) y TD-A1 (administracion),
grupo electrogeno con ATS, UPS para combustible/control y para
POS-CCTV-comunicaciones, alumbrado normal y de emergencia, tomacorrientes,
bombas sumergibles de tanque, cabezas electronicas de surtidores, control
ATG, compresor de aire, bombas de agua y de efluentes, puesta a tierra,
equipotencialidad, proteccion contra el rayo y propuesta academica de
clasificacion de areas peligrosas.

No incluye factibilidad de concesionaria, calculo definitivo de
cortocircuito/selectividad, obra civil, instalaciones de GLP/GNV, proyecto
sanitario ni contra incendio completo, ni habilitacion profesional. Estos
componentes requieren documentos o especialidades que no fueron entregados.

\section{Suministro y continuidad}
Se adopta 3~\(\times\)~380/220~V, 60~Hz, 3F+N+PE, suministro directo en baja
tension con capacidad de servicio propuesta de \ServCap. La empresa
distribuidora se asume como Electro Puno por la ubicacion del grifo, y la
propuesta queda condicionada a su factibilidad y a la corriente de
cortocircuito disponible en el punto de entrega. El neutro y el conductor de
proteccion se mantienen separados aguas abajo del origen (TN-S).

Para continuidad se propone un grupo \GrupoModelo~de~\GrupoNom~kVA standby con
ATS de cuatro polos y UPS senoidales para combustible/control y para
POS-CCTV-comunicaciones. El respaldo es un criterio de ingenieria del
anteproyecto, no una obligacion general atribuida sin matiz a todos los
grifos.
